\documentclass[11pt,a4paper]{article}
\usepackage[T1]{fontenc}
\usepackage[utf8]{inputenc}
\usepackage{lmodern}
\usepackage[margin=25mm]{geometry}
\usepackage{amsmath,amssymb,amsthm,mathtools}
\usepackage{booktabs,longtable,array}
\usepackage{enumitem}
\usepackage{microtype}
\usepackage{xcolor}
\usepackage{hyperref}
\hypersetup{colorlinks=true,linkcolor=blue!55!black,urlcolor=blue!55!black,
  citecolor=blue!55!black,pdftitle={SIMBA: Mathematical Foundations and Implementation Contract},
  pdfauthor={Ajit Kumar}}
\setlist{itemsep=3pt,topsep=5pt}
\setlength{\emergencystretch}{2em}
\newtheorem{theorem}{Theorem}[section]
\newtheorem{proposition}[theorem]{Proposition}
\newtheorem{corollary}[theorem]{Corollary}
\newtheorem{definition}[theorem]{Definition}
\theoremstyle{remark}
\newtheorem{remark}[theorem]{Remark}
\newcommand{\R}{\mathbb{R}}
\newcommand{\E}{\mathbb{E}}
\newcommand{\Var}{\operatorname{Var}}
\newcommand{\Cov}{\operatorname{Cov}}
\newcommand{\diag}{\operatorname{diag}}
\newcommand{\supp}{\operatorname{supp}}
\newcommand{\ind}{\mathbf{1}}
\newcommand{\inc}{\Delta_{\mathrm{inc}}}
\newcommand{\calI}{\mathcal{I}}
\newcommand{\calB}{\mathcal{B}_{\mathrm{tr}}}
\newcommand{\code}[1]{\texttt{\small #1}}
\title{SIMBA: Mathematical Foundations\\and Implementation Contract}
\author{Ajit Kumar}
\date{15 October 2026}

\begin{document}
\maketitle
\begin{abstract}
SIMBA is a bundle-completion procedure combining a pairwise collaborative
backbone, historical joint-incidence surplus, redundancy-weighted quadratic
seed concordance, and exposure-modulated Gaussian shrinkage. This document
specifies the operators implemented in the standalone repository and derives
their mathematical properties. It proves a separation between pairwise
projection and triple incidence, characterizes concordance through elementary
symmetric polynomials and Rao diversity, derives the normal--normal posterior,
and states the exact conditions under which inverse-frequency variance
components are identifiable. It also distinguishes algebraic identities,
model-conditional statistical results, implementation policies, and empirical
retrieval claims. The constituent results on hypergraphs, diversity, and
Gaussian conjugacy are established mathematics; their use here does not
constitute a claim to have invented those primitives.
\end{abstract}
\tableofcontents
\newpage

\section{Purpose, scope, and types of claims}
This is the detailed theoretical companion to \code{simba/operators.py}.
The shorter paper methodology is in \code{paper/method.tex}. The discussion
applies to the current operators, rather than to every historical SIMBA
experiment. Dataset-specific evaluators add feature channels and select
coefficients; they do not all instantiate an identical scoring formula.

Four kinds of statements must be kept separate:
\begin{enumerate}
  \item \textbf{Algebraic facts:} identities and operator properties that
  hold for the specified inputs, independent of a data-generating model.
  \item \textbf{Model-conditional results:} posterior or moment-estimation
  statements that require explicit distributional or conditional-moment
  assumptions.
  \item \textbf{Software policies:} choices such as positive variance floors,
  clipping, validation grids, and disabling unidentified shrinkage.
  \item \textbf{Empirical claims:} improvements in Recall@20 or NDCG@20,
  which require an evaluation protocol and measured evidence.
\end{enumerate}
For example, joint incidence can contain information absent from a pairwise
projection. This does not prove that a particular incidence bonus improves
ranking. Similarly, a posterior mean minimizes squared-error Bayes risk under
its assumed model, but that result is not a Recall@20 guarantee.

\section{Task, notation, and information boundary}
\subsection{Bundle completion}
Let $\calI=\{1,\ldots,n\}$ be a fixed item catalogue and
$\calB=\{b_1,\ldots,b_m\}$ the training bundles. Each bundle is a set;
repeated item identifiers within a bundle are removed by the graph
constructor. A query consists of an observed seed set $S\subset\calI$.
The candidate domain is $\calI\setminus S$. At evaluation time, a disjoint
held-out set $H$ supplies the relevant completion items.

Define the binary incidence matrix and degrees by
\begin{equation}
 B_{bi}=\ind\{i\in b\},\qquad
 f_i=\sum_b B_{bi},\qquad d_b=\sum_i B_{bi}.
 \label{eq:incidence}
\end{equation}
The seed indicator is $q_S\in\{0,1\}^n$; a batch of $Q$ queries forms
$J=[q_{S_1},\ldots,q_{S_Q}]\in\{0,1\}^{n\times Q}$.
Unless explicitly stated otherwise, all graph results below use binary $B$.
The optional weighted constructor is treated separately in
Section~\ref{sec:weighted}.

\begin{center}
\begin{tabular}{@{}ll@{}}
\toprule
Symbol & Meaning \\
\midrule
$a_i=(f_i+5)^{-\alpha}$ & Seed exposure attenuation \\
$v_b=(d_b+1)^{-\nu}$ & Bundle-volume attenuation \\
$r_b=(d_b+1)^{-1}$ & Smoothing bundle normalization \\
$\beta\geq0$ & Joint-incidence coefficient \\
$\tau>0$ & Concordance denominator regularizer \\
$w_{cs}\geq0$ & Candidate support from seed $s$ \\
$G_S$, $K_S$ & Seed similarity and concordance weighting matrices \\
$Y_{cS}$, $\mu_{cS}$ & Fused observation score and scaled prior mean \\
$\sigma_0^2$, $\sigma_\epsilon^2$ & Prior and exposure-noise variance components \\
$\kappa=\sigma_\epsilon^2/\sigma_0^2$ & Exposure shrinkage scale \\
\bottomrule
\end{tabular}
\end{center}
The graph defaults are $\alpha=\nu=1/2$ and $\beta=1.5$.
The code calls $\nu$ \code{gamma\_vol}; it is unrelated to the concordance
kernel parameter $\gamma$ used later.

\subsection{What each split may supply}
Training bundles determine collaborative incidence and exposure. Catalogue
metadata and frozen content features describe items, including items unseen
in training bundles. Validation queries determine score scaling, variance
components, and any selected fusion or shrinkage parameters. Test targets
determine evaluation metrics only. Catalogue-level text preprocessing may
use all item descriptions; that is a transductive catalogue assumption,
not evidence that the corresponding held-out bundle memberships are known.

The benchmark loaders also have dataset-specific auxiliary interactions,
filters, and query caps. Those choices are documented in
\code{docs/REPRODUCIBILITY.md}. A mathematical property of the shared operator
does not certify a loader's experimental protocol.

\section{Pairwise induction and propagation}
\subsection{Additive pairwise support}
Set $D_a=\diag(a)$ and $D_v=\diag(v)$. Before diagonal removal, the
pairwise score matrix is
\begin{equation}
 W^{\mathrm{full}}=B^\top D_v B D_a,
 \qquad W=W^{\mathrm{full}}-\diag(\diag(W^{\mathrm{full}})).
 \label{eq:pairgraph}
\end{equation}
Thus $W_{cs}=a_s\sum_{b:\,c,s\in b}v_b$ for $c\neq s$. The attenuation
is on the seed side, so $W$ need not be symmetric. The additive candidate
score is
\begin{equation}
 A^{\mathrm{pair}}_c(S)=\sum_{s\in S}W_{cs}.
 \label{eq:pairscore}
\end{equation}
\code{score\_pairwise\_and\_diff2} computes $W^{\mathrm{full}}J$ in
factored form. For $c\notin S$, this is exactly~\eqref{eq:pairscore}.
Seed rows can retain self-contributions in the intermediate array; the
evaluators exclude seed items from the final candidate ranking. Those
intermediate values can still enter propagated channels.

\subsection{Symmetric smoothing is not a Markov transition}
The binary-incidence implementation of \code{propagate} applies
\begin{equation}
 P_{\mathrm{sym}}=D_a B^\top D_r B D_a-D_h,
 \quad h_i=a_i^2\sum_b r_b B_{bi},\quad D_r=\diag(r).
 \label{eq:psym}
\end{equation}
Its diagonal is zero; for $i\neq j$,
$(P_{\mathrm{sym}})_{ij}=a_i a_j\sum_b r_b B_{bi}B_{bj}\geq0$.
The matrix is symmetric, but its rows are not generally normalized.
Removing the diagonal also need not preserve positive semidefiniteness.
Consequently, symmetry alone gives neither a probability transition nor
a contraction in a chosen norm. ``Diffusion'' here denotes graph smoothing.
The function returns $P_{\mathrm{sym}}X$; naming a channel ``diff2'' does
not mean that this function itself computes $P_{\mathrm{sym}}^2X$.

\subsection{The separately exposed random-walk operator}
For $f_i>0$ and $d_b>0$, let
\begin{equation}
 T_0=D_{1/f} B^\top D_{1/d} B,\qquad
 \bar T=T_0-\diag(\diag(T_0)),\qquad
 \rho_i=\sum_j\bar T_{ij}.
\end{equation}
Zero-degree reciprocals are defined to be zero. Define
\begin{equation}
 (P_{\mathrm{rw}})_{ij}=\begin{cases}
   \bar T_{ij}/\rho_i,&\rho_i>0,\\
   0,&\rho_i=0.
 \end{cases}
 \label{eq:prw}
\end{equation}
\begin{proposition}[Active-row stochasticity]
$P_{\mathrm{rw}}$ has nonnegative entries and zero diagonal. Every row with
$\rho_i>0$ sums to one; every other row is zero. On active rows,
$(P_{\mathrm{rw}}x)_i$ is a convex combination of neighbouring values.
\end{proposition}
\begin{proof}
All factors of $T_0$ are nonnegative. Diagonal deletion preserves off-diagonal
nonnegativity. Dividing a positive row sum by itself gives one, and a
zero row sum in a nonnegative matrix implies an all-zero row. The convex
combination statement follows directly from these weights.
\end{proof}
A warm item occurring only in singleton bundles can also have a zero row.
Warm exposure alone is therefore insufficient for the active-row condition.
The evaluators do not all use $P_{\mathrm{rw}}$; some use~\eqref{eq:psym}
or a graph mixture specific to their dataset. These operators are related
to established hypergraph constructions~\cite{zhou2006}, but the precise
normalizations above specify the implementation contract.

\section{Historical joint-incidence surplus}
\subsection{Definition and exact matrix form}
For a query $S$, define
\begin{equation}
 o_b(S)=\sum_{s\in S}B_{bs}=|b\cap S|,\qquad
 u_b(S)=\sum_{s\in S}a_s B_{bs}.
\end{equation}
The pure surplus is
\begin{equation}
 \inc(c\mid S)=\beta\sum_b B_{bc}v_b u_b(S)
                      \max\{0,o_b(S)-1\}.
 \label{eq:delta}
\end{equation}
For a batch, with $\odot$ denoting elementwise multiplication,
\begin{equation}
 D_{\mathrm{inc}}=\beta B^\top
 \left[(D_v B D_aJ)\odot(BJ-\mathbf 1)_+\right].
 \label{eq:delta-matrix}
\end{equation}
\code{score\_simplicial\_incidence\_only} returns this array and its
propagated version. It does not add~\eqref{eq:pairscore}.

\begin{proposition}[Nonnegativity and exact activation]
Assume $a_s>0$, $v_b>0$, and $\beta>0$. For $c\notin S$,
$\inc(c\mid S)>0$ if and only if a training bundle contains $c$ and at
least two distinct seeds. In particular, $|S|\leq1$ implies $\inc=0$.
With $\beta=0$, the surplus is identically zero.
\end{proposition}
\begin{proof}
Each summand in~\eqref{eq:delta} is nonnegative. It is strictly positive
exactly when $B_{bc}=1$ and $o_b\geq2$, since these conditions also imply
$u_b>0$. A finite nonnegative sum is positive exactly when one summand is.
\end{proof}

\subsection{What order of interaction is actually represented?}
\begin{proposition}[Triple-incidence representation]
Let $T_{cst}=\sum_b v_b B_{bc}B_{bs}B_{bt}$ for distinct $c,s,t$.
For candidates outside $S$,
\begin{equation}
 \inc(c\mid S)=\beta\sum_{\{s,t\}\subset S}(a_s+a_t)T_{cst}.
 \label{eq:triple}
\end{equation}
Thus the direct surplus is a query-quadratic aggregation of weighted
three-item incidences, even when training hyperedges contain more than
three items.
\end{proposition}
\begin{proof}
For a particular bundle, write $x_s=B_{bs}\in\{0,1\}$.
Since $x_s^2=x_s$,
\begin{align}
 u_b(o_b-1)
 &=\sum_s a_s x_s\left(\sum_t x_t-1\right)\\
 &=\sum_{s\neq t}a_s x_sx_t
 =\sum_{\{s,t\}\subset S}(a_s+a_t)x_sx_t.
\end{align}
If $o_b=0$, both $u_b$ and the expression are zero; if $o_b\geq1$, the
clamp in~\eqref{eq:delta} does not change the product. Substituting this
identity into~\eqref{eq:delta} gives~\eqref{eq:triple}.
\end{proof}
This distinction matters. The operator accesses genuinely joint information
absent from a pairwise graph, but it is not an unrestricted estimator of
all $k$-way interactions for every $k$. Two hypergraphs with matching $T$
and matching normalization statistics produce the same direct surplus.
There is no boundary operator, orientation, Hodge Laplacian, or homology
calculation in this function. ``Simplicial'' names its joint-incidence
construction; it does not establish a topological learning guarantee.

\subsection{Separation from identical pairwise projections}
Consider the following unweighted hypergraphs on items $0,\ldots,5$:
\begin{align}
 \mathcal H_1&=\{012,034,135,245\},\\
 \mathcal H_2&=\{013,024,125,345\}.
 \label{eq:pasch}
\end{align}
Here $012$ denotes the set $\{0,1,2\}$. Each item has degree two, all
bundles have size three, and the off-diagonal pairs in both systems are
\begin{equation*}
 \{01,02,12,03,04,34,13,15,35,24,25,45\}.
\end{equation*}
Consequently, the full $B^\top B$, including its diagonal, agrees between
the systems. So do $W$, $P_{\mathrm{sym}}$, and $P_{\mathrm{rw}}$.

\begin{theorem}[Pairwise information is insufficient]
No function of the common pairwise projection and its shared degree/size
statistics can reproduce~\eqref{eq:delta} for every hypergraph.
\end{theorem}
\begin{proof}
For $S=\{0,1\}$ and $c=2$, $\mathcal H_1$ has the contributing bundle
$012$, giving $\inc=\beta v_{012}(a_0+a_1)>0$. In $\mathcal H_2$, no
bundle contains $0,1,2$, so $\inc=0$. A deterministic function of the
identical pairwise inputs must give the same output in both cases and
therefore cannot equal both surplus values.
\end{proof}
This is an expressivity separation, not an identification theorem for a
whole hypergraph, and not a statistical advantage on every dataset.
The matched sizes avoid confounding the comparison with volume weights.

\subsection{Activation rate, propagation, and retrieval}
Define the query event
\begin{equation}
 E(S)=\{\exists b\in\calB: |b\cap S|\geq2\}.
\end{equation}
For a random query, let $A(S)$ mean that at least one candidate outside $S$
has a nonzero direct surplus. Then
\begin{equation}
 \Pr(A(S))\leq\Pr(E(S)).
 \label{eq:activation}
\end{equation}
The inequality can be strict: an overlapping bundle may contain seeds only.
If $E(S)$ is false, the entire surplus vector is zero and linear propagation
also returns zero. If it is true, propagation can deliver signal to candidates
that do not themselves share the original hyperedge.

\begin{corollary}[An opportunity bound on a bounded metric]
Suppose two scoring pipelines differ only by this surplus and its deterministic
propagation, with identical remaining parameters and tie handling. For a
per-query metric $M\in[0,1]$,
\begin{equation}
 \left|\E[M_{\mathrm{aug}}-M_{\mathrm{base}}]\right|
 \leq\Pr(E(S)).
\end{equation}
\end{corollary}
\begin{proof}
Outside $E(S)$ the two scores and rankings agree, so the difference is zero.
Inside the event the metric difference has absolute value at most one.
Apply the triangle inequality and take expectations.
\end{proof}
This loose bound applies only to a controlled comparison with everything
else fixed. Retuning global coefficients after adding the channel can
change inactive-query scores and invalidates that argument. The bound
does not determine the sign of improvement on active queries.

\section{Quadratic concordance from candidate--seed supports}
\subsection{Nonnegative support construction}
For each candidate and seed, the Rao operator uses
\begin{equation}
 w_{cs}=W_{cs}+\eta\max\{0,x_c^\top x_s\},\qquad \eta\geq0,
 \label{eq:support}
\end{equation}
when dense features are supplied. Otherwise the second term is absent.
Theoretical cosine bounds below assume unit-norm feature rows; zero rows
may be retained. The function does not normalize arbitrary inputs itself.
Write $A_c=\sum_{s\in S}w_{cs}$. If $A_c>0$, set $p_{cs}=w_{cs}/A_c$;
when $A_c=0$, define all concordance scores to be zero.

\subsection{Elementary symmetric polynomial identity}
Let $e_2(w_c)=\sum_{s<t}w_{cs}w_{ct}$. The unweighted operator is
\begin{equation}
 R_0(c\mid S;\tau)=\frac{2e_2(w_c)}{A_c+\tau}
 =\frac{A_c^2-\sum_s w_{cs}^2}{A_c+\tau}.
 \label{eq:r0}
\end{equation}
The standalone \code{score\_concordance\_e2} function uses $w_{cs}=W_{cs}$
without the dense addition. Equation~\eqref{eq:r0} also describes the
unweighted kernel case of the general Rao formula for any supplied
nonnegative supports, including~\eqref{eq:support}.
\begin{proposition}[Mass--diversity factorization]
For $A_c>0$,
\begin{equation}
 R_0=\frac{A_c^2}{A_c+\tau}
       \left(1-\sum_s p_{cs}^2\right).
 \label{eq:gini}
\end{equation}
If exactly $k_+>0$ seeds have positive support, then
\begin{equation}
 0\leq R_0\leq\frac{A_c^2}{A_c+\tau}\left(1-\frac1{k_+}\right),
 \label{eq:r0bound}
\end{equation}
with equality at the upper bound for equal support on those $k_+$ seeds.
\end{proposition}
\begin{proof}
Expanding $(\sum_s w_{cs})^2$ yields the sum of squares plus $2e_2$.
Substitution of $w_{cs}=A_cp_{cs}$ proves~\eqref{eq:gini}.
Cauchy--Schwarz gives $1=(\sum p_{cs})^2\leq k_+\sum p_{cs}^2$.
The lower bound follows from nonnegative cross-products, and equal shares
attain the upper bound.
\end{proof}
At fixed total mass, concordance rewards dispersed support. Without fixing
mass, it trades diversity against magnitude. The code clamps a small
negative value of $A_c^2-\sum_s w_{cs}^2$ to zero to guard cancellation.

\subsection{Relation to products and its exact limit}
For $m_S=|S|$,
\begin{equation}
 \prod_{s\in S}(\tau+w_{cs})
 =\sum_{j=0}^{m_S}\tau^{m_S-j}e_j(w_c),\qquad e_0=1.
 \label{eq:product}
\end{equation}
This identity explains how $e_2$ occurs as one term in a product expansion.
It does not make $R_0$ equal to a product-of-experts likelihood: other
terms are omitted and the denominator is different. Strict conjunction
would depend on all seeds, for example through $e_{m_S}=\prod_s w_{cs}$.
For equal total mass three,
\begin{align*}
 w^{(1)}&=(1.5,1.5,0),& e_2(w^{(1)})&=2.25,&e_3(w^{(1)})&=0,\\
 w^{(2)}&=(2.1,0.45,0.45),& e_2(w^{(2)})&=2.0925,&e_3(w^{(2)})&=0.42525.
\end{align*}
Concordance prefers the first, while the strict product prefers the second.
The use of second-order terms resembles established bilinear models such
as factorization machines~\cite{rendle2010}; the support normalization
and seed kernel specify the particular score used here.

\section{Redundancy weighting through Rao diversity}
\subsection{Similarity and weighting matrices}
With category labels $\ell_s$ and normalized features, the default combined
seed similarity is
\begin{equation}
 G_{st}=\tfrac12\ind\{\ell_s=\ell_t\}
          +\tfrac12\max\{0,x_s^\top x_t\},\quad s\neq t,
 \qquad G_{ss}=1.
 \label{eq:similarity}
\end{equation}
With only one source, that source determines $G$. With neither, the code
uses $G=I$. For $s\neq t$, its general kernel is
\begin{equation}
 K^{(\gamma)}_{st}=\begin{cases}
   \max\{0,1+\gamma G_{st}\},&\gamma<0,\\
   1+\gamma G_{st},&\gamma\geq0,
 \end{cases}
 \qquad K^{(\gamma)}_{ss}=0.
 \label{eq:kernel}
\end{equation}
Complementarity mode has $\gamma=-1$, giving $K_{st}=1-G_{st}$.
The music evaluator instead uses $\gamma=1$, which favors agreement
between similar seeds. That is a coherence weighting, not the same
redundancy penalty. In either case, missing or inaccurate metadata changes
what ``different'' or ``similar'' means.

\subsection{Quadratic diversity and implementation}
The score is
\begin{equation}
 R_K(c\mid S;\tau)=\frac{w_c^\top K_Sw_c}{A_c+\tau}
  =\frac{A_c^2}{A_c+\tau}\,p_c^\top K_Sp_c.
 \label{eq:rao}
\end{equation}
For a dissimilarity matrix $K$, $p^\top Kp$ is Rao's quadratic
diversity~\cite{rao1982}. It is also the expected dissimilarity between
two independent draws of a seed index from $p$. The independence here is
a mathematical device for interpreting the double sum; it does not assert
independence of the observed seed evidence. The sum counts both orders of
each unordered seed pair, consistently with $2e_2$ in~\eqref{eq:r0}.

\begin{proposition}[Bounds and vanishing conditions]
If $K$ is symmetric, zero-diagonal, and $0\leq K_{st}\leq K_{\max}$,
then
\begin{equation}
 0\leq R_K\leq K_{\max}\frac{A_c^2}{A_c+\tau}
                  \left(1-\sum_s p_{cs}^2\right)
 \leq K_{\max} A_c.
 \label{eq:raobound}
\end{equation}
The score is zero if every pair with positive support has $K_{st}=0$.
It is positive exactly when at least one pair has positive support on
both seeds and a strictly positive kernel entry.
\end{proposition}
\begin{proof}
Expand $w^\top Kw=\sum_{s\neq t}K_{st}w_sw_t$. All terms are
nonnegative. Bounding each kernel entry by $K_{\max}$ and applying
\eqref{eq:gini} gives~\eqref{eq:raobound}. The positivity condition follows
from the same finite nonnegative sum.
\end{proof}
Complementarity mode has $K_{\max}=1$. For normalized features and
$\gamma=1$, $K_{\max}=2$. These bounds concern raw scores, before
catalogue z-normalization. Neither positive semidefiniteness nor a metric
triangle inequality is needed for these bounds.

\subsection{Redundant and diverse seeds}
If all seeds have $G_{st}=1$, complementarity-mode concordance is zero.
If all off-diagonal similarities are zero, $R_K=R_0$. This formalizes
the difference between redundant and distinct support sources, conditional
on the metadata proxy. Different categories are not by themselves proof
of functional compatibility; dissimilar items can also be incompatible.

The operator is invariant to permutation of the same seed set when the
corresponding rows and columns of $K$ are permuted together. It is not
generally invariant to adding a distinct item identifier with identical
metadata to an existing seed. Such an addition changes $A_c$, the support
distribution, and its cross-terms with other seed groups. Exact
duplication-invariance would require a separately specified group-level
aggregation; it is not a property claimed here.

Nor is $R_K$ coordinatewise monotone for every admissible $K$. For example,
take three seeds with $K_{12}=K_{21}=1$ and all other entries zero.
At $w=(1,1,t)$,
\begin{equation}
 R_K(w;\tau)=\frac{2}{2+t+\tau},
\end{equation}
which decreases with $t$. A new support source that contributes no
dissimilarity-weighted agreement can dilute the normalized score.

\subsection{Scaling law and exact finite-scale correction}
For a fixed seed kernel and any $\lambda>0$,
\begin{equation}
 R_K(\lambda w;\lambda\tau)=\lambda R_K(w;\tau).
 \label{eq:homogeneous}
\end{equation}
This follows because the numerator scales quadratically and the denominator
linearly. For fixed $\tau$, writing $Q_K=w^\top Kw$ instead yields
\begin{equation}
 R_K(\lambda w;\tau)=\frac{\lambda^2 Q_K}{\lambda A+\tau}.
 \label{eq:fixedtau}
\end{equation}
For $Q_K>0$, the ratio to $\lambda R_K(w;\tau)$ is
\begin{equation}
 \frac{R_K(\lambda w;\tau)}{\lambda R_K(w;\tau)}
 =\frac{\lambda(A+\tau)}{\lambda A+\tau}.
\end{equation}
Thus exact degree-one homogeneity requires joint scaling. The large-support
approximation replaces $A+\tau$ by $A$; its relative shortfall from
$Q_K/A$ is exactly $\tau/(A+\tau)$. Merely letting $\lambda$ grow with
fixed original $A$ does not prove equality to $\lambda R_K(w;\tau)$.
If support is rescaled by changing underlying model parameters, the dense
and graph terms in~\eqref{eq:support} must both scale to invoke
\eqref{eq:homogeneous}.

\subsection{Information separation from incidence}
Rao concordance is a function of candidate--seed supports and seed
similarities. If those inputs agree, its output agrees. In particular,
the two hypergraphs in~\eqref{eq:pasch}, with matched metadata, yield
identical concordance but different incidence surplus. Conversely,
concordance can be positive when its supports come from distinct pairwise
bundles and no single bundle contains two seeds. The channels therefore
access different information. This does not imply statistical independence,
uncorrelated outputs, or orthogonality of fitted estimators.

\section{Fusion, metadata priors, and interpretation of the target}
\subsection{Catalogue normalization and fused observations}
For a channel $g_k(c,S)$, define column-wise catalogue normalization
\begin{equation}
 z_k(c,S)=\frac{g_k(c,S)-\bar g_k(S)}{s_k(S)+\varepsilon_z},
 \qquad \bar g_k(S)=\frac1n\sum_i g_k(i,S).
 \label{eq:zscore}
\end{equation}
The repository uses $\varepsilon_z=10^{-9}$. NumPy uses its default
population standard deviation, while the Torch implementation uses its
default sample correction. This small convention difference is part of
the implementation, not a new statistical estimator.
Typically,
\begin{equation}
 Y_{cS}=\sum_k\omega_k z_k(c,S),\qquad
 P_{cS}=\sum_{k\in\mathcal P}\xi_k z_k(c,S),\qquad
 \mu_{cS}=s_{\mathrm{val}}P_{cS}.
 \label{eq:fusion}
\end{equation}
The selected $\mathcal P$ need not contain only exogenous metadata:
the controlled retail prior includes bridge, text, auxiliary interaction,
and Rao terms. $Y$ and $P$ can therefore share channels and training data.
The term ``prior'' describes its role in shrinkage, not a guarantee that
it is independent evidence. The moment model below must be assessed for
the actual fused scores, including this dependence.

\subsection{The latent target is a model, not an observed label}
Let $\theta_{cS}\in\R$ denote latent compatibility on the standardized
score scale. It is not defined here as a probability, a calibrated
log-odds, or a statistically identified log-lift. The score $Y$ is a
constructed collaborative/content estimator, not a direct repeated
measurement of an observed $\theta$. The Gaussian hierarchy supplies a
working model connecting the two. Algebraically exact posterior updating
under that model is a narrower statement than correct real-world
calibration of $\theta$.

\section{Exposure-modulated Gaussian posterior}
\subsection{Assumptions}
For fixed $S,c$, positive exposure $f_c$, and a fitted prior mean $\mu_{cS}$,
assume
\begin{align}
 \theta_{cS}\mid\mu_{cS},f_c
   &\sim\mathcal N(\mu_{cS},\sigma_0^2),\\
 Y_{cS}\mid\theta_{cS},\mu_{cS},f_c
   &\sim\mathcal N(\theta_{cS},\sigma_\epsilon^2/f_c),
 \label{eq:hierarchy}
\end{align}
with strictly positive variances. The second line specifies a zero-mean
likelihood error conditional on the latent compatibility and prior mean.
It does not follow merely from positive co-occurrence counts. Frequency
is treated as an exposure proxy, although real histories can have unequal
opportunity, temporal drift, and dependent observations.

\subsection{Derivation by completing the square}
Suppress the indices and put $v_f=\sigma_\epsilon^2/f$. The posterior
log-density, up to constants independent of $\theta$, is
\begin{align}
 \log p(\theta\mid Y,\mu,f)
 &=-\frac12\left[\frac{(\theta-\mu)^2}{\sigma_0^2}
                 +\frac{(Y-\theta)^2}{v_f}\right]+C\\
 &=-\frac12\left[(\sigma_0^{-2}+v_f^{-1})\theta^2
       -2(\mu\sigma_0^{-2}+Yv_f^{-1})\theta\right]+C'.
\end{align}
Therefore the posterior is Gaussian with
\begin{align}
 V_f&=\left(\frac1{\sigma_0^2}+\frac f{\sigma_\epsilon^2}\right)^{-1},
 \label{eq:postvar}\\
 \widehat\theta
 &=V_f\left(\frac\mu{\sigma_0^2}+\frac{fY}{\sigma_\epsilon^2}\right)
 =\frac{fY+\kappa\mu}{f+\kappa}
 =(1-u_f)Y+u_f\mu,\quad
 u_f=\frac\kappa{f+\kappa},\quad
 \kappa=\frac{\sigma_\epsilon^2}{\sigma_0^2}.
 \label{eq:posterior}
\end{align}
This is the literal update in \code{conjugate\_posterior} and
\code{apply\_eb\_shrinkage}; no positive-part residual is added.

\begin{proposition}[Convexity, exposure limits, and sensitivity]
For $f\geq0$ and $\kappa>0$, $0<u_f\leq1$ and the posterior mean lies
between $Y$ and $\mu$. Moreover,
\begin{equation}
 \frac{\partial u_f}{\partial f}=-\frac\kappa{(f+\kappa)^2},\qquad
 \frac{\partial\widehat\theta}{\partial Y}=\frac f{f+\kappa},\qquad
 \frac{\partial\widehat\theta}{\partial\mu}=\frac\kappa{f+\kappa}.
\end{equation}
At $f=0$ the zero-likelihood-precision limit is $\widehat\theta=\mu$;
as $f\to\infty$, $\widehat\theta\to Y$ for fixed $Y,\mu,\kappa$.
\end{proposition}
The $f=0$ case is a limit and a software convention; the positive-exposure
Gaussian likelihood in~\eqref{eq:hierarchy} is not a proper finite-variance
observation distribution at $f=0$. In contrast, \code{kappa\_f=0}
explicitly disables shrinkage and returns $Y$ even on cold items. It
should not be interpreted by substituting $f=\kappa=0$ in the fraction.

For $f=(0,0,9)$, $Y=(-2,-2,4)$, $\mu=(-1,1,1)$ and $\kappa=4$,
\begin{equation}
 \widehat\theta=(-1,1,40/13)\approx(-1,1,3.076923).
\end{equation}
Coldness alone does not impose a positive score: a negative cold prior
remains negative. Whether that score enters a top-$K$ list still depends
on every competing candidate.

\subsection{A conditional squared-error guarantee, with its limits}
\begin{theorem}[Bayes risk under the specified hierarchy]
For known variances and a correctly specified hierarchy, the posterior
mean minimizes conditional expected squared error among all measurable
estimators based on $Y,\mu,f$. Its integrated risk for fixed $f>0,\mu$ is
\begin{equation}
 \E[(\widehat\theta-\theta)^2\mid f,\mu]
 =\frac{\sigma_0^2\sigma_\epsilon^2}{f\sigma_0^2+\sigma_\epsilon^2}
 \leq\min\{\sigma_0^2,\sigma_\epsilon^2/f\}.
 \label{eq:risk}
\end{equation}
\end{theorem}
\begin{proof}
For any $a=a(Y,\mu,f)$, the conditional squared-error decomposition is
\begin{equation}
 \E[(\theta-a)^2\mid Y,\mu,f]
 =\Var(\theta\mid Y,\mu,f)
   +(a-\E[\theta\mid Y,\mu,f])^2.
\end{equation}
The second term is minimized at the posterior mean. The posterior variance
is the constant $V_f$ in~\eqref{eq:postvar}, so integrating over $Y$
gives~\eqref{eq:risk}. Its denominator is at least either individual
precision contribution, giving both inequalities.
\end{proof}
This is an established property of posterior means, not a new SIMBA
optimality theorem. Estimated or validation-selected hyperparameters,
biased $Y$, prior misspecification, and correlated prior/observation errors
can invalidate the comparison. Lower latent squared error also does not
imply better retrieval of held-out binary memberships. The implementation
returns a plug-in mean; it does not return uncertainty intervals accounting
for fitted hyperparameters or a verified coverage guarantee.

\section{Variance-component estimation and identifiability}
\subsection{From the hierarchy to a conditional moment}
Write $\theta-\mu=\zeta$ and $Y-\theta=\epsilon$, so $Y-\mu=\zeta+\epsilon$.
Under~\eqref{eq:hierarchy},
\begin{equation}
 \E[(Y-\mu)^2\mid f,\mu]
 =\sigma_0^2+\sigma_\epsilon^2/f.
 \label{eq:moment}
\end{equation}
The cross-term vanishes because
$\E[\epsilon\mid\zeta,f,\mu]=0$.
Gaussianity is sufficient for the posterior derivation, but not necessary
for this moment: zero means, the stated variances, and a zero cross-moment
are enough. Conversely, the moment alone does not imply a Gaussian
posterior.

With $Q_v$ validation queries, the implemented item-level response is
\begin{equation}
 v_c=\frac1{Q_v}\sum_{S\in\mathcal Q_v}(Y_{cS}-\mu_{cS})^2,
 \qquad x_c=1/f_c,\quad c\in\mathcal W=\{c:f_c>0\}.
 \label{eq:responses}
\end{equation}
The method regresses $v_c$ on $x_c$; it does not regress the variance of
$Y$ alone, nor the spread of the prior means. Each warm item receives
one OLS observation regardless of its exposure.

\subsection{Raw OLS formulas and their proof conditions}
Let $N_w=|\mathcal W|$, $D_x=\sum_{c\in\mathcal W}(x_c-\bar x)^2$.
When $N_w\geq2$ and $D_x>0$, the raw estimates are
\begin{align}
 \widetilde\sigma_\epsilon^2
   &=\frac{\sum_c(x_c-\bar x)(v_c-\bar v)}{D_x},\\
 \widetilde\sigma_0^2&=\bar v-\widetilde\sigma_\epsilon^2\bar x.
 \label{eq:ols}
\end{align}
They solve the unweighted least-squares normal equations for an intercept
and slope. Calling these values variance components describes their
intended target; finite-sample values may be negative.

\begin{proposition}[Conditional unbiasedness of raw components]
For a fixed identifiable design vector $x$, suppose
$v_c=\sigma_0^2+\sigma_\epsilon^2x_c+e_c$ and
$\E[e_c\mid x]=0$ for each warm item. Then the raw slope and intercept
in~\eqref{eq:ols} are conditionally unbiased.
\end{proposition}
\begin{proof}
Substitution into the slope gives
\begin{equation}
 \widetilde\sigma_\epsilon^2
 =\sigma_\epsilon^2+\frac{\sum_c(x_c-\bar x)e_c}{D_x}.
\end{equation}
Conditional expectation of the second term is zero. The intercept error
is $\bar e-\bar x(\widetilde\sigma_\epsilon^2-\sigma_\epsilon^2)$,
which also has conditional expectation zero.
\end{proof}
Independence across items and homoskedastic errors are not needed for this
expectation calculation. They do matter for usual OLS standard errors,
efficiency, and concentration. Shared queries and graph structure can
correlate residuals, so treating $N_wQ_v$ cells as independent observations
would be unjustified. The routine does not estimate regression standard
errors.

There is an additional fitted-score caveat: $Y$ and $\mu$ are constructed
and their scale or coefficients can be estimated on these same validation
queries. The conditional-mean-zero premise is not automatically preserved
after that fitting and selection. The raw formulas are OLS; unconditional
finite-sample unbiasedness of the complete fitted SIMBA pipeline has not
been established. A separate calibration split or cross-fitting could
make such an analysis cleaner, but is not implemented by this document.

\subsection{Identifiability and the numerical decision rule}
\begin{proposition}[Moment identifiability]
The intercept and slope are separately identifiable from the linear
moment equation exactly when the design matrix with columns $1,x$ has
rank two, equivalently when at least two distinct warm inverse frequencies
are observed. With a single common $x_0$, only
$\sigma_0^2+x_0\sigma_\epsilon^2$ is identified.
\end{proposition}
\begin{proof}
The normal matrix has determinant
\begin{equation}
 \det\begin{pmatrix}N_w&\sum x_c\\\sum x_c&\sum x_c^2\end{pmatrix}
 =N_wD_x.
\end{equation}
It is invertible exactly when $D_x>0$. For constant $x_0$, every pair of
components with the same intercept-plus-$x_0$-times-slope has the same
fitted conditional moment.
\end{proof}
The code uses the stronger numerical requirement $N_w\geq2$ and
$D_x>10^{-12}$, testing a \emph{sum} of squared centered inverse
frequencies, not a sample variance threshold. It otherwise returns
\code{identifiable=False}, null raw and constrained components, and
\code{kappa\_f=0.0}. All production evaluator grids then restrict their
candidate shrinkage values to zero. Non-identification is not an estimate
that the true noise component equals zero; disabling shrinkage is the
chosen fallback policy.

\subsection{Positive truncation and ratio clipping}
For identifiable designs, the returned plug-in components are
\begin{align}
 \widehat\sigma_0^2&=\max\{\widetilde\sigma_0^2,10^{-4}\},&
 \widehat\sigma_\epsilon^2&=\max\{\widetilde\sigma_\epsilon^2,10^{-4}\},\\
 \widehat\kappa&=\operatorname{clip}
   (\widehat\sigma_\epsilon^2/\widehat\sigma_0^2,0.05,50).
 \label{eq:clip}
\end{align}
The raw outputs are unrounded; constrained outputs are reported separately.
The nonlinear transformations in~\eqref{eq:clip} are generally biased,
and even an unclipped ratio of unbiased components need not be unbiased.
These componentwise floors are not the joint solution of a constrained
nonnegative least-squares problem. ``Constrained plug-in'' refers to
admissible returned values, not to an optimization problem that the code
does not solve. Some evaluators subsequently choose a different positive
$\kappa$ from a validation grid; that is a tuned posterior-shaped score,
not necessarily the direct moment-ratio estimate.

\subsection{Prior-shift invariance and misspecification}
For any array $h_{cS}$, replacing $(Y,\mu)$ by $(Y+h,\mu+h)$ leaves
$v_c$ and the variance estimates unchanged. This exact residual invariance
explains why a simulated change in the spread of prior means need not
change $\widehat\kappa$ when the same prior is added to $Y$. It does not
prove that estimating a different prior on real data is harmless.

For example, if the residual has conditional bias $b(f)$ and prior/noise
cross-covariance $q(f)$, its second moment can instead contain
\begin{equation}
 \E[(Y-\mu)^2\mid f]
 =\sigma_0^2+\sigma_\epsilon^2/f+b(f)^2+2q(f),
\end{equation}
under an otherwise comparable decomposition. Frequency-dependent prior
error or correlated reused channels can then be mistaken for exposure
noise. Diagnostics should examine residual-moment curves against $1/f$,
negative raw components, clipping frequency, and sensitivity to the prior.

\section{Validation-frozen inference and a batch-invariance result}
\subsection{Calibration contract}
A typical prior scale is
\begin{equation}
 s_{\mathrm{val}}=\frac{\operatorname{std}(Y_{\mathrm{val}})}
 {\operatorname{std}(P_{\mathrm{val}})+10^{-9}},
 \label{eq:scale}
\end{equation}
where the standard deviations pool the validation score array. Individual
evaluators differ slightly in stabilizers and whether scale one is also
considered. Every selected scale is frozen before test inference. The
same applies to variance ratios, channel weights, encoder parameters,
and any backward-elimination decisions. A moment fit uses validation
scores without requiring held-out target labels; validation metric
selection can use validation targets. Neither operation is performed
on the test-query score batch in the current pipelines.

\begin{theorem}[Score independence from unrelated query columns]
Fix the catalogue, trained operators, metadata, and calibration parameters.
Suppose each channel for query $S$ is computed only from $S$ and those
fixed objects; any z-normalization is across catalogue items within its
own query column. Then appending unrelated query columns to a batch
does not change $\widehat\theta_{cS}$ in exact arithmetic.
\end{theorem}
\begin{proof}
The raw channel column for $S$ is unchanged by hypothesis. Its mean and
standard deviation in~\eqref{eq:zscore} use that column only, so the
normalized column is unchanged. Fixed linear fusion gives unchanged
$Y_{cS}$ and $P_{cS}$. Frozen scaling gives unchanged $\mu_{cS}$, and
\eqref{eq:posterior} uses only these values and fixed $f_c,\kappa$.
\end{proof}
Re-estimating~\eqref{eq:scale} on each test batch would violate the frozen
parameter premise. Floating-point kernels may introduce tiny layout-dependent
differences, and evaluator tie perturbations can affect exact tied ranks;
the proposition concerns the deterministic mathematical score. It does
not assert invariance to changing the catalogue, because catalogue means,
standard deviations, and available candidates can then change.

\section{Auxiliary channels and boundaries of the theory}
\subsection{Ontology and slot saturation}
The ontology operator forms category co-occurrence statistics and
positive-PMI-like off-diagonal weights, with optional soft clusters and
within-type channels. \code{score\_slot\_saturation} adds a penalty based
on query seeds sharing a candidate's category and empirical category
duplication rates. Despite its name, this is an additive pairwise diagonal
correction. It is not a distinct higher-order joint-incidence estimator.
Its interpretation depends on the category system and on whether a
bundle is expected to fill distinct slots or contain repeated types.

\subsection{Spectral features and content bridges}
The graph operator also constructs positive-PMI co-occurrence features
and a truncated-SVD embedding. Content bridges mix frozen features with
one or more graph propagation steps. Their scores supply additional
channels and help score training-cold items with catalogue features.
These operations do not make cold items collaboratively observed, and
their feature quality is not guaranteed by the incidence or posterior
proofs. A content-neighbour matrix can be row-normalized on populated
rows without making every graph smoother in the system stochastic.

\subsection{Set-attention encoder}
The full evaluators additionally train a small neural set encoder on
training examples, using transformed item features, attention pooling,
and a contrastive objective with sampled negatives. Its output is fused
as another channel. A set sum or attention aggregation is permutation
invariant when attention weights are computed equivariantly and no
order-dependent input is introduced. That structural observation is
not a convergence or generalization theorem for the trained encoder.
The offline demo and controlled retail ablation do not train this model;
their execution does not verify the full encoder's optimization.

\subsection{Pruning and retrieval policies}
\code{prune\_weights\_val} performs validation-tolerance backward
elimination of optional channels. It is not a minimum-description-length
objective, and no global optimum over channel subsets is claimed.
The optional tail/cold slot reservation is a post-hoc retrieval policy,
separate from the normal--normal posterior. If enabled, the resulting
ranking is not solely the ordering of~\eqref{eq:posterior}. Reports
should state whether this policy was used.

\section{Controlled ablations: what is and is not fixed}
In the retail ablation, let $L$ denote the collaborative pairwise mixture,
$P_g$ the fixed dataset-specific graph propagation, and $B_r,T,U$ the
bridge, text, and auxiliary interaction scores. Writing $z$ for catalogue
normalization, the implemented ladder is
\begin{align}
 F_0&=z(L)+0.8z(P_gL)+0.8z(B_r)+0.5z(T)+0.5z(U),\\
 F_1&=z(L+\inc)+0.8z(P_g(L+\inc))
                 +0.8z(B_r)+0.5z(T)+0.5z(U),\\
 F_2&=F_0+0.35z(R_K),\\
 F_3&=F_1+0.35z(R_K),\\
 F_4&=(1-u)F_3+u\mu.
 \label{eq:ablations}
\end{align}
The prior before frozen scaling is
$P=0.8z(B_r)+0.5z(T)+0.5z(U)+0.35z(R_K)$.
The graph, data split, raw backbone channels, and numerical coefficients
are shared within a seed. Only validation scores fit the final prior
scale and moment ratio.

Because $z(L+\inc)\neq z(L)+z(\inc)$ in general, adding incidence also
changes the normalization of the combined collaborative channel. The
comparison measures the implemented incidence-plus-normalization
intervention, not a fixed standardized baseline plus a raw surplus.
Likewise $F_4-F_3=u(\mu-F_3)$ both attenuates and adds evidence; it is not
a pure cold-item bonus. The $F_1-F_0$ and $F_2-F_0$ comparisons isolate
the two additions, while $F_3$ checks their combination and $F_4$ adds
posterior calibration. This is a controlled ladder with two individual
mechanism branches, not a complete factorial study of all interactions.

For a held-out set $H$, Recall@$K$ is
\begin{equation}
 \operatorname{Recall@}K(S,H)
 =\frac{|\operatorname{TopK}(\calI\setminus S)\cap H|}{|H|}.
\end{equation}
Split-seed means and sample standard deviations describe split variation.
A positive mean difference alone does not establish significance or
causality. Paired per-query or per-split comparisons, exposure strata,
activation-conditioned metrics, and explicit tie handling strengthen
interpretation. Current evidence and its provenance remain in
\code{docs/RESULTS.md}; this theory document adds no new benchmark results.

\section{Computation, memory, and weighted-input caveats}
\subsection{Operator costs}
Let $L_B=\operatorname{nnz}(B)$, $L_W=\operatorname{nnz}(W)$, and
$k=|S|$. Factored incidence for $Q$ dense output columns takes on the
order of $L_BQ$ sparse--dense work plus bundle/query intermediates and
an $nQ$ output. It avoids storing an $n^3$ triple tensor even though
\eqref{eq:triple} characterizes its information. Computing graph diffusion
in factored form also uses order $L_BQ$ work.

Unweighted concordance can use $WJ$ and $W^{\odot2}J$, with sparse
products rather than explicit enumeration of seed pairs. In the general
Rao implementation, each query materializes an $n\times k$ support
array and multiplies it by a $k\times k$ kernel: the dense quadratic
work is $O(nk^2)$, with support memory $O(nk)$, plus output storage and
any dense feature products. A dense feature width $d$ adds $O(ndk)$
work for candidate--seed dot products. Full-catalogue score banks can
therefore require several GiB even when a neural encoder has few parameters.

Preprocessing costs also matter: forming $B^\top B$ can generate up to
order $\sum_b d_b^2$ pair contributions before aggregation. Sparse input
does not imply that its projection remains sparse. SVD, auxiliary graphs,
clustering, and model training add costs beyond the scoring formulas.
No sublinear full-catalogue retrieval guarantee is claimed.

\subsection{Optional weighted constructor}\label{sec:weighted}
The constructor accepts bundle weights. Let $H$ be binary incidence,
$t_b\geq0$ the supplied weight, and $\widetilde B=D_tH$ the stored
nonbinary matrix. Then exposure is $f_i=\sum_b t_bH_{bi}$, bundle sizes
come from $H$, and a bilinear projection uses
\begin{equation}
 \widetilde B^\top D_v\widetilde B
 =H^\top\diag(t_b^2v_b)H.
\end{equation}
Weights are squared because they enter both factors. The surplus also
contains $t_b^2$, while overlap counts use $H$. It is not a once-weighted
hyperedge sum unless all weights are one or their interpretation is
adjusted accordingly.

There is a further binary-specific diagonal correction in
\code{propagate}. For weighted input the implemented matrix is
\begin{equation}
 \widetilde P=D_a\widetilde B^\top D_r\widetilde B D_a
       -\diag\left(a^2\odot(\widetilde B^\top r)\right),
\end{equation}
whose diagonal entry is
\begin{equation}
 \widetilde P_{ii}=a_i^2\sum_b r_b(t_b^2-t_b)H_{bi}.
\end{equation}
Thus the zero-diagonal and nonnegative-smoother statements for binary
$B$ do not automatically extend to arbitrary bundle weights. The shipped
evaluators instantiate the constructor without this optional weighting.
These caveats specify the existing function rather than silently proving
a property for a different weighted operator.

\section{Verification map and remaining scientific questions}
\subsection{What the existing suite checks}
The mathematical audit checks active-row stochasticity, pure incidence
activation, separation of a joint bundle from a pairwise triangle, the
unweighted concordance identity, redundancy examples, numerical posterior
values, simulated variance-component recovery, equal-frequency
non-identification, frozen-scale batch invariance, validation pruning,
and production evaluator compliance with the non-identification policy.
The standalone suite checks packaging, data utilities, the offline demo,
and fail-closed evaluation orchestration.

Unit tests instantiate examples and catch implementation regressions;
they do not prove all quantified statements in this document. In
particular, the fully degree-matched pairwise separation example in
\eqref{eq:pasch} is justified by the explicit enumeration above, rather
than by assuming the simpler triangle test matches every diagonal statistic.

\subsection{What remains empirical or unresolved}
The following are scientific questions rather than consequences of the
preceding algebra:
\begin{enumerate}
 \item Does training exposure approximate inverse observation variance
 after normalization and feature fusion? Is $1/f$ better than alternatives?
 \item Does the metadata-derived $K$ distinguish redundant from
 complementary constraints in the target application?
 \item Does incidence help on active queries, especially after matching
 popularity, query size, and candidate frequency?
 \item Are posterior gains robust to prior quality, channel reuse,
 split selection, and variance clipping? How often are raw components
 negative or numerically unidentified?
 \item Do gains persist against strong tuned methods under the same
 candidate catalogue, data access, split, and computational budget?
 \item Can a meaningful generative target or retrieval guarantee be
 established beyond a working standardized-score hierarchy?
\end{enumerate}
The present algorithm combines distinct mechanisms and exposes testable
predictions about their behavior. Its mathematical guarantees are those
proved above, under their stated domains. Novelty of the full method and
its publication value additionally require related-work comparison and
empirical evidence; renaming established estimators does not supply that
evidence.

\appendix
\section{Reference inference procedure}
\begin{enumerate}
 \item Construct train-only incidence, degrees, collaborative graphs,
 and any allowed auxiliary interaction channels. Prepare the fixed
 catalogue features under the declared transductive or inductive protocol.
 \item For validation queries, compute pairwise, incidence, concordance,
 content, and optional learned scores. Normalize each query column across
 the fixed catalogue; construct $Y_{\mathrm{val}}$ and $P_{\mathrm{val}}$.
 \item Fit the validation prior scale, compute residual moments, and
 return both raw and admissible variance components. If unidentified,
 restrict the shrinkage grid to zero. Complete any validation selection.
 \item Freeze all fitted objects and parameters. For each test query,
 compute its own score column and apply the chosen posterior formula.
 Do not refit scales or variance components on test-input batches.
 \item Exclude seeds, apply the declared tie convention and any explicitly
 enabled post-hoc policy, then select the top $K$ eligible candidates.
 \item Compute metrics from held-out test targets only after scoring.
 Record split seeds, data paths/provenance, parameters, and policy settings.
\end{enumerate}

\section{Source map and build instructions}
The authoritative detailed source is \code{docs/THEORY.tex}. A navigation
page is retained at \code{docs/THEORY.md}; a compiled copy is supplied as
\code{docs/THEORY.pdf}. From the repository root, build with a standard
LaTeX installation:
\begin{verbatim}
mkdir -p artifacts/theory
for pass in 1 2 3; do
  pdflatex -halt-on-error -interaction=nonstopmode \
    -output-directory=artifacts/theory docs/THEORY.tex || exit 1
done
\end{verbatim}
Repeated passes resolve the table of contents and internal references.
No BibTeX step or external figure is required. Generated files stay under
the ignored \code{artifacts/} directory; copy the final PDF to
\code{docs/THEORY.pdf} when refreshing its checked-in rendering.
Core implementation locations are:
\begin{center}
\begin{tabular}{@{}p{0.39\linewidth}p{0.54\linewidth}@{}}
\toprule
File & Role \\
\midrule
\code{simba/operators.py} & Shared graph, incidence, concordance, ontology,
and shrinkage functions \\
\code{simba/ablation\_study.py} & Controlled retail scoring ladder and
validation-frozen moment calibration \\
\code{simba/*\_eval.py} & Dataset-specific feature banks, learned channels,
validation selection, and retrieval \\
\code{tests/test\_audit.py} & Mathematical and evaluator-policy examples \\
\code{tests/test\_standalone.py} & Installation/data/driver checks \\
\bottomrule
\end{tabular}
\end{center}

\begin{thebibliography}{9}
\bibitem{zhou2006}
D. Zhou, J. Huang, and B. Sch\"olkopf.
Learning with Hypergraphs: Clustering, Classification, and Embedding.
\emph{Advances in Neural Information Processing Systems 19}, 2006.
\href{https://papers.nips.cc/paper_files/paper/2006/hash/dff8e9c2ac33381546d96deea9922999-Abstract.html}
{NeurIPS proceedings, original paper}.

\bibitem{rao1982}
C. R. Rao. Diversity and dissimilarity coefficients: a unified approach.
\emph{Theoretical Population Biology}, 21(1):24--43, 1982.
\url{https://doi.org/10.1016/0040-5809(82)90004-1}.

\bibitem{rendle2010}
S. Rendle. Factorization Machines.
\emph{IEEE International Conference on Data Mining}, 2010.
\url{https://doi.org/10.1109/ICDM.2010.127}.
\end{thebibliography}
\end{document}
